% A Practical Preimage Attack on 19-Round SHA-256
% via a Global sigma0-Difference Table
%
% Compile: pdflatex paper_r19_final.tex  (twice for cross-references)

\documentclass[11pt,a4paper]{article}

\usepackage{amsmath,amssymb,amsthm}
\usepackage{booktabs}
\usepackage{xcolor}
\usepackage[hidelinks]{hyperref}
\hypersetup{pdftitle={A Practical Preimage Attack on the 19-Round SHA-256 Compression Function via a Global sigma0-Difference Table},pdfauthor={Kameldip Singh Basra},pdfsubject={Cryptanalysis of step-reduced SHA-256},pdfkeywords={SHA-256, preimage attack, step-reduced, message schedule, table lookup, GPU cryptanalysis}}
\usepackage{geometry}
\geometry{margin=1in}
\usepackage[protrusion=true,expansion=false]{microtype}

\newtheorem{theorem}{Theorem}
\newtheorem{corollary}{Corollary}
\newtheorem{definition}{Definition}
\newtheorem{lemma}{Lemma}
\newtheorem{proposition}{Proposition}
\newtheorem{remark}{Remark}

\newcommand{\Sig}[1]{\Sigma_{#1}}
\newcommand{\sig}[1]{\sigma_{#1}}
\newcommand{\Maj}{\mathrm{Maj}}
\newcommand{\Ch}{\mathrm{Ch}}
\newcommand{\ROTR}[2]{\mathrm{ROTR}^{#1}(#2)}
\newcommand{\SHR}[2]{\mathrm{SHR}^{#1}(#2)}
\newcommand{\mmod}{\bmod\,2^{32}}

\begin{document}

\title{\textbf{A Practical Preimage Attack on the\\
       19-Round SHA-256 Compression Function\\
       via a Global \texorpdfstring{$\sigma_0$}{sigma0}-Difference Table}}

\author{Kameldip Singh Basra \\[4pt]
        \small Independent Researcher \\
        \small \texttt{kameldipbasra@gmail.com} \\
        \small ORCID: \href{https://orcid.org/0009-0001-0977-4614}{0009-0001-0977-4614}}

\date{September 2026}

\maketitle

%-----------------------------------------------------------------------
\begin{abstract}
We present a practical preimage attack on the compression function of
SHA-256 reduced to 19 of its 64 rounds, in the standard setting: the
chaining value is the fixed initial value, the target is an arbitrary
256-bit digest, and the attacker chooses the 512-bit message block.  The
attack combines a backward chain that fixes eight state words from the
target, a precomputed $16$\,GB table inverting $u\mapsto\sigma_0(u)-u$, and
a two-stage fixed-point iteration over the three schedule-consistency
constraints.  Each attempt sweeps one $32$-bit variable, costs $O(2^{32})$
table lookups, and runs in about $25$\,s on one NVIDIA H100 with a single seed chain ($89$\,s with the four chains used in the campaign).  Seven preimages were found and independently verified.  Six are of digests of random messages: three from dedicated runs that succeeded in their 4th, 21st and 35th contexts, and three from a preregistered campaign of 240 attempts that gives a direct estimate of the per-attempt success rate, $\hat\mu = 3/240 =
1.25\times10^{-2}$ (95\% Poisson interval $[2.6\times10^{-3},\,
3.7\times10^{-2}]$; all three events fell in the screened arm and the
control arm gave $0/120$).  Attempts are far from interchangeable: in a
paired campaign with matched exposure, a screening pass costing $1/4096$ of
an attempt raised the yield of the intermediate 16-bit filter by $3.41\times$ (95\% CI $[2.13, 5.82]$, target-blocked bootstrap; counted per converged trajectory, with identity-level deduplication bounding it within $[2.42\times, 3.44\times]$), helping on 31 of 40 targets and hurting on 9; its effect on the final success rate is not established.

SAT-based cryptanalysis inverts 17 and 18 rounds in seconds on one core,
and Zaikin has recently inverted the 19-round compression function for a
fixed all-ones digest with a parameterized Cube-and-Conquer solver, in
$18.5$ hours on 192 CPU cores.  Priority at 19 rounds is his.  The
contribution here is a different route to the same round count, algebraic
rather than search-based, that reaches an arbitrary digest in about two hours of one GPU at the pooled campaign rate (one attempt in eighty; the unscreened control arm gave $0/120$, so the figure for a uniformly random context is nearer three hours), did so for his target in 84 attempts ($123$ minutes; the seventh preimage, and a second preimage of that digest) after a first run of 14 attempts had produced a match that a since-fixed bug discarded, and comes with complete artefacts.  A companion note explains why the fourth schedule constraint costs this construction a further factor of $2^{32}$ at 20 rounds; a note added at the end points to a later chosen-context variant that pays that factor and reaches 20 rounds.  Full SHA-256 is not affected by any of
these results.

\medskip
\noindent\textbf{Keywords:} SHA-256, preimage attack, step-reduced,
message schedule, table lookup, GPU cryptanalysis.
\end{abstract}

%-----------------------------------------------------------------------
\section{Introduction}
\label{sec:intro}
%-----------------------------------------------------------------------

SHA-256~\cite{FIPS180} is the most widely deployed member of the SHA-2
family.  Its compression function has 64 rounds, and the security margin of
the full function is usually assessed through attacks on step-reduced
variants.

\paragraph{Attack model.}
Let $\mathrm{compress}_R$ denote the SHA-256 compression function
restricted to its first $R$ rounds and let $\mathrm{IV}$ be the standard
initial value.  Given a 256-bit digest $h$, a \emph{preimage} is a 512-bit
block $(W_0,\ldots,W_{15})$ with
\[
  h \;=\; \mathrm{IV} + \mathrm{compress}_R(\mathrm{IV},\,W_0,\ldots,W_{15}) .
\]
This is the standard preimage setting for the compression function: the
chaining value is fixed, nothing about any original message is known, and
only the digest is given.  It is not a pseudo-preimage or free-start
setting, in which the attacker would also choose the chaining value.  The
16 words of the block are unconstrained, so the result concerns the
compression function and not the length-padded hash.  For a single-block
55-byte message the padding fixes $W_{14}$, $W_{15}$ and the low byte of
$W_{13}$; those three conditions can be met by choosing three of the
attack's context words in closed form (\texttt{padded\_preimage.py} in the
artefacts), but a 167-attempt run produced no padded preimage where one to two were expected at the projected $6\times10^{-3}$ and at the pooled campaign rate ($P(0)$ between $0.12$ and $0.29$), so the run neither confirms nor excludes the padded variant and we claim only the compression-function result.

\paragraph{Prior work.}
Theoretical preimage attacks on step-reduced SHA-256 reach far beyond 19
rounds, but at costs within a few bits of the generic $2^{256}$: Isobe and
Shibutani~\cite{IsobeShibutani2009} give a one-block preimage attack on 24
steps at $2^{240}$; Aoki et al.~\cite{AokiGuo2009} reach 43 steps at
$2^{254.9}$ by meet-in-the-middle; Khovratovich, Rechberger and
Savelieva~\cite{KRS2012} reach 45 steps at $2^{255.5}$ with bicliques; Guo, Li, Liu, Wang and Zhang~\cite{GuoLiLiuWangZhang2026} have since extended the meet-in-the-middle line to 46 steps at $2^{254.5}$ and 47 steps at $2^{255.6}$.  None of these has been executed in full.  Collision attacks are a different problem and are far more advanced: Li, Liu and Wang~\cite{LiLiuWang2024} lowered the 31-step collision attack on SHA-256 to $2^{49.8}$ and Li et al.~\cite{LiLiuWangDongSun2024} then computed a practical 31-step collision, both building on local-collision techniques~\cite{SanadhyaSarkar2008,MendelNadSchlaffer2013}; they do not yield preimages.

The practical preimage frontier belongs to SAT-based cryptanalysis, which Mironov and Zhang~\cite{MironovZhang2006} introduced for hash functions in the collision setting.  Davydov, Pikh\-tov\-ni\-kov,
Kir\-ya\-no\-va and Zaikin~\cite{Davydov2025} find preimages of the 17- and 18-round
compression function with Kissat under six CNF encodings (CBMC, Transalg, and four adder variants of SAT-encoding), report 19 rounds as unsolved under every one of them, and reach 19 rounds only in a weakened
form in which at most 31 of the 256 output bits are constrained, and they also report 19-round collisions, a different problem; their published benchmark instances target the all-zero digest.  For SHA-256 both literatures use ``round'' and ``step'' for the same 64 applications of the step function, so no prior result changes its round count under the other convention.
Zaikin~\cite{Zaikin2026} then constructs intermediate inverse problems
between consecutive round counts, parameterizes a CDCL solver on them, and
reports that the resulting parallel Cube-and-Conquer solver inverted the
29-step MD5, 24-round SHA-1 and 19-round SHA-256 compression functions for
the first time.  For SHA-256 the target is one fixed digest, the all-ones
string, with all 256 bits fixed, in the same model as here: compression
function, standard IV, no padding, feed-forward kept.  The 19-round
instance itself was not solved within seven days on 192 cores; the preimage
was obtained from the intermediate problem between rounds 18 and 19 whose
solutions are 19-round preimages with probability $1/8$, in
$18$\,h\,$33$\,min on 192 CPU cores (two 96-core AMD EPYC 9654), about
$148$ core-days, and it verifies against the reference implementation.
That article was submitted in March 2025 and appeared in January 2026.
Priority at 19 rounds is Zaikin's and we claim none.

\paragraph{This work.}
We give an algebraic, table-based preimage attack on the 19-round
compression function.  Its cost profile is explicit and low: after a
one-time $16$\,GB precomputation, each attempt is a single $2^{32}$ sweep
on one GPU, $25$ to $89$ seconds on an NVIDIA H100 depending on the number
of seed chains, and in the campaign of \S\ref{sec:screening} (four chains)
about one attempt in eighty succeeded when the two arms are pooled ($3/240$; 95\% interval $[2.6\times10^{-3},\,3.7\times10^{-2}]$), which at $89$\,s per attempt is about two GPU-hours per preimage.  That pooled figure is the headline cost quoted in this paper, but it mixes a screened arm ($3/120$, about one GPU-hour) with a control arm of unscreened contexts that gave $0/120$; the screened arm was the top $30\%$ of a screening pool, so the population-weighted point estimate for a uniformly random context is $0.3\times3/120+0.7\times0/120=7.5\times10^{-3}$ per attempt, about $3.3$ GPU-hours.  Three events cannot separate these figures: the pooled interval spans a factor of 14 and contains all of them.  The ingredients are:
\begin{enumerate}
  \item a backward chain that recovers the eight state words
    $a_{11},\ldots,a_{18}$ from the digest alone (\S\ref{sec:structure});
  \item a global table inverting $u\mapsto\sigma_0(u)-u$, built once and
    valid for every target and every attempt (\S\ref{sec:table});
  \item the observation that in each of the three schedule constraints one
    unknown state word appears with opposite signs in two places, so that
    the constraint reduces to a single table lookup once the remaining
    unknowns are provisionally fixed (\S\ref{sec:attack}); the provisional
    assignment is closed by a fixed-point iteration and tested exactly by a
    32-bit consistency check;
  \item seven independently verified preimages: six of digests of random
    messages, three from dedicated runs and three from a preregistered
    240-attempt campaign that also measures the success rate directly, and
    one of the all-ones digest that Zaikin attacked, found in 84 attempts
    (\S\ref{sec:experiments}, \S\ref{sec:comparison});
  \item a measured, preregistered account of how much a cheap screening
    pass improves the intermediate yield, and of what it does not establish
    (\S\ref{sec:screening}).
\end{enumerate}
A companion note~\cite{Basra2026R20} explains why the same construction does not reach 20 rounds at the 19-round cost (see the note added at the end of this paper).  Appendix~\ref{app:ninestep} records an exact
cancellation identity in the message schedule found along the way,
together with the reason the search experiments built on it are not
claimed as results.

%-----------------------------------------------------------------------
\section{SHA-256 Preliminaries}
\label{sec:prelim}
%-----------------------------------------------------------------------

\paragraph{Notation.}
All arithmetic is modulo $2^{32}$ unless stated otherwise.
We use $\oplus$, $\mathbin{\&}$, $\overline{x}$ for bitwise XOR, AND, NOT.
Rotations and shifts of a 32-bit word $x$:
\[
  \ROTR{n}{x} = (x \gg n) \mid (x \ll (32-n)),
  \qquad
  \SHR{n}{x} = x \gg n.
\]

\paragraph{Round functions.}
\begin{align*}
  \Sig{0}(a) &= \ROTR{2}{a} \oplus \ROTR{13}{a} \oplus \ROTR{22}{a}, \\
  \Sig{1}(e) &= \ROTR{6}{e} \oplus \ROTR{11}{e} \oplus \ROTR{25}{e}, \\
  \sig{0}(x) &= \ROTR{7}{x}  \oplus \ROTR{18}{x} \oplus \SHR{3}{x},  \\
  \sig{1}(x) &= \ROTR{17}{x} \oplus \ROTR{19}{x} \oplus \SHR{10}{x}, \\
  \Ch(e,f,g) &= (e \mathbin{\&} f) \oplus (\overline{e} \mathbin{\&} g), \\
  \Maj(a,b,c) &= (a \mathbin{\&} b) \oplus (a \mathbin{\&} c)
                 \oplus (b \mathbin{\&} c).
\end{align*}

\paragraph{State labels.}
We write the 8-word compression state after round $r$ as
$(a_r,a_{r-1},a_{r-2},a_{r-3},\allowbreak e_r,e_{r-1},e_{r-2},e_{r-3})$,
with IV boundary values $a_{-4},\ldots,a_{-1}$
and $e_{-4},\ldots,e_{-1}$:
\[
  a_{-1}{=}\mathrm{IV}[0],\; a_{-2}{=}\mathrm{IV}[1],\;
  a_{-3}{=}\mathrm{IV}[2],\; a_{-4}{=}\mathrm{IV}[3],\;
  e_{-1}{=}\mathrm{IV}[4],\; e_{-2}{=}\mathrm{IV}[5],\;
  e_{-3}{=}\mathrm{IV}[6],\; e_{-4}{=}\mathrm{IV}[7].
\]

\paragraph{Round equations.}
\begin{align}
  T_1^{(r)} &= \Sig{1}(e_{r-1}) + \Ch(e_{r-1},e_{r-2},e_{r-3})
               + e_{r-4} + K_r + W_r, \label{eq:T1}\\
  T_2^{(r)} &= \Sig{0}(a_{r-1}) + \Maj(a_{r-1},a_{r-2},a_{r-3}),
               \label{eq:T2}\\
  a_r       &= T_1^{(r)} + T_2^{(r)}, \label{eq:a} \\
  e_r       &= a_{r-4} + T_1^{(r)}. \label{eq:e}
\end{align}

Note that $e_{r-4}$ in~\eqref{eq:T1} is the \emph{h-register} word
(boundary value $e_{-4}=\mathrm{IV}[7]$), while $a_{r-4}$ in~\eqref{eq:e}
is the \emph{d-register} word (boundary value $a_{-4}=\mathrm{IV}[3]$);
these are distinct IV words and distinct chain values for all $r\ge 4$.

From~\eqref{eq:a}--\eqref{eq:e}: $T_1^{(r)} = a_r - T_2^{(r)}$,
and rearranging~\eqref{eq:T1}:
\begin{equation}
  W_r = T_1^{(r)} - \Sig{1}(e_{r-1}) - \Ch(e_{r-1},e_{r-2},e_{r-3})
        - e_{r-4} - K_r.
  \label{eq:W-round-inv}
\end{equation}

\paragraph{Message schedule.}
For $r \geq 16$ the message schedule expands:
\begin{equation}
  W_r = \sig{1}(W_{r-2}) + W_{r-7} + \sig{0}(W_{r-15}) + W_{r-16}.
  \label{eq:schedule}
\end{equation}
For $r < 16$ the word $W_r$ is an independent input to the compression
function; it is what we seek.

\paragraph{Hash output.}
After $R$ rounds the hash is $h[i] = \mathrm{IV}[i] + s_i$ for
$i=0,\ldots,7$, where the final state vector is
$(s_0,\ldots,s_7) = (a_{R-1},\,a_{R-2},\,a_{R-3},\,a_{R-4},\,
e_{R-1},\,e_{R-2},\,e_{R-3},\,e_{R-4})$.

%-----------------------------------------------------------------------
\section{Structural Analysis}
\label{sec:structure}
%-----------------------------------------------------------------------

\subsection{Backward Chain}

From the $R$-round hash we recover the final state directly:
\[
  a_{R-1}, a_{R-2}, a_{R-3}, a_{R-4},\;
  e_{R-1}, e_{R-2}, e_{R-3}, e_{R-4}.
\]
We then extend backward using
$T_2^{(r)} = \Sig{0}(a_{r-1}) + \Maj(a_{r-1},a_{r-2},a_{r-3})$,
$T_1^{(r)} = a_r - T_2^{(r)}$,
$a_{r-4} = e_r - T_1^{(r)}$,
iterating for $r = R-1, R-2, R-3, R-4$ to recover
$a_{R-5},\ldots,a_{R-8}$.

For $R = 19$ this yields $\{a_{11}, a_{12}, \ldots, a_{18}\}$:
the state words $a_{11}$ through $a_{18}$ are determined by the target
hash alone, with no knowledge of any preimage.

\subsection{Free Parameters and Schedule Constraints}

At $R=19$ the unconstrained parameters are $a_0,\ldots,a_{10}$
(11 words).  Of these, $a_4,\ldots,a_{10}$ are chosen uniformly at
random per attempt (\emph{context}), forming 7 attacker-controlled
free words.  The remaining $a_0,\ldots,a_3$ are the targets of the
attack phases.

The schedule~\eqref{eq:schedule} first activates at $r=16$, yielding
three consistency constraints at $R=19$:
\begin{align}
  W_{16} &= \sig{1}(W_{14}) + W_9 + \sig{0}(W_1) + W_0, \tag{C0}\\
  W_{17} &= \sig{1}(W_{15}) + W_{10} + \sig{0}(W_2) + W_1, \tag{C1}\\
  W_{18} &= \sig{1}(W_{16}) + W_{11} + \sig{0}(W_3) + W_2. \tag{C2}
\end{align}
Here $W_{16}, W_{17}, W_{18}$ are recovered from the target alone via the
backward chain and~\eqref{eq:W-round-inv}, using the known
$a_{11},\ldots,a_{18}$ together with the context values
$a_4,\ldots,a_{10}$ (which supply $e_{12},e_{13},e_{14}$ through $a_8,a_9,a_{10}$; together with $e_{15},e_{16},e_{17}$ from the chain alone these are the $e$-words that~\eqref{eq:W-round-inv} needs at $r=16,17,18$).  Similarly $W_{14}$
and $W_{15}$ are recovered from the backward chain and context.

%-----------------------------------------------------------------------
\section{The \texorpdfstring{$\sigma_0$}{sigma0}-Difference Table}
\label{sec:table}
%-----------------------------------------------------------------------

\begin{definition}
\label{def:table}
The \emph{$\sigma_0$-difference table} $\mathcal{T}$ is an array
of $2^{32}$ 32-bit entries:
\[
  \mathcal{T}[v] = u
\]
where $u$ is a representative satisfying $\sig{0}(u) - u \equiv v
\pmod{2^{32}}$, or a sentinel $\perp$ if no such $u$ exists.
When multiple $u$ share the same $v$, the largest such $u$ is stored; this is the policy of the published table and of every run reported here, and \S\ref{sec:complexity} (E3) tests the alternative.  On disk the table is an \texttt{int32} array with sentinel $-1$; the single value $v=\mathtt{0x20000000}$, whose only root is $u=\mathtt{0xFFFFFFFF}$, then coincides with the sentinel, one entry in $2.7\times10^{9}$.  The GPU build uses the \texttt{int32} minimum as sentinel and has no such collision.
\end{definition}

\paragraph{Construction.}
$\mathcal{T}$ is built in a single pass: for each $u \in [0, 2^{32})$,
compute $v = (\sig{0}(u) - u) \bmod 2^{32}$ and write $u$ to
$\mathcal{T}[v]$.  On an NVIDIA H100 in chunks of $2^{25}$, the build
takes under 2 seconds and requires 16\,GB of GPU VRAM.

\paragraph{Coverage.}
The function $f(u) = \sig{0}(u) - u \bmod 2^{32}$ is not injective.
Empirically, $2{,}721{,}603{,}628$ of the $2^{32}$ possible values
$v$ have at least one preimage, corresponding to \textbf{63.37\%}
coverage.

\paragraph{Usage.}
Given a target value $v^*$, the query $\mathcal{T}[v^*]$ returns in
$O(1)$ a word $u$ satisfying $\sig{0}(u) - u = v^*$, or $\perp$
(probability $\approx 36.6\%$).

%-----------------------------------------------------------------------
\section{The 19-Round Preimage Attack}
\label{sec:attack}
%-----------------------------------------------------------------------

\subsection{Overview}

The attack operates in three phases per context:

\begin{enumerate}
  \item[\textbf{C0.}] Sweep $a_0 \in [0, 2^{32})$.  A table lookup
    recovers $W_1$ (and hence $a_1$) for about $63.4\%$ of $a_0$ values (the image fraction of the table).
  \item[\textbf{C1/C2.}] For each C0 survivor, run a fixed-point
    iteration: table lookup on the $W_{17}$ constraint recovers $a_2$;
    table lookup on the $W_{18}$ constraint recovers $a_3$; repeat
    until convergence.
  \item[\textbf{W9 filter.}] Test the consistency condition
    $\varepsilon(a_2,a_3)=0$ of Proposition~\ref{lem:w9-diff} exactly; its
    pass rate is measured in \S\ref{sec:complexity}.  Candidates that pass
    are verified by forward evaluation.
\end{enumerate}

The key that makes all three phases work is a common algebraic pattern:
in each constraint, the target unknown $a_i$ appears with \emph{opposite
signs} in two different terms, causing exact cancellation in the
$\sigma_0$-difference lookup target.

\subsection{Context Precomputation}

Fix a target hash $h$.  Run the backward chain to obtain
$a_{11},\ldots,a_{18}$.  Choose uniformly random context
$a_4,\ldots,a_{10}$.  Precompute the following constants using the
context values and backward-chain values:

\begin{align}
  T_2^{\mathrm{IV}} &= \Sig{0}(a_{-1}) + \Maj(a_{-1},a_{-2},a_{-3}),
  \label{eq:T2iv} \\
  C_0 &= {-T_2^{\mathrm{IV}} - e_{-4} - \Sig{1}(e_{-1})
          - \Ch(e_{-1},e_{-2},e_{-3}) - K_0}, \label{eq:C0const} \\
  C_{e_0} &= {a_{-4} - T_2^{\mathrm{IV}}}.  \label{eq:Ce0}
\end{align}

From these: $W_0 = a_0 + C_0$ and $e_0 = a_0 + C_{e_0}$, both
linear in $a_0$.

For the context rounds $r = 7,\ldots,11$, compute
$T_2^{(r)}, T_1^{(r)}, e_r$ from~\eqref{eq:T1}--\eqref{eq:e} using
the context $a_4,\ldots,a_{10}$ together with $a_{11}$ from the backward
chain.
Define:
\begin{align}
  W_9^{\mathrm{base}} &= T_1^{(9)} - \Sig{1}(e_8) - K_9, \label{eq:W9base}\\
  W_{11}^{\mathrm{base}} &= T_1^{(11)} - T_1^{(7)} - \Sig{1}(e_{10})
                             - \Ch(e_{10},e_9,e_8) - K_{11}, \label{eq:W11base}\\
  \mathit{CONST}_{10} &= T_1^{(10)} - \Sig{1}(e_9) - K_{10}. \label{eq:C10}
\end{align}

Also recover $W_{14}, W_{15}, W_{16}^{\mathrm{bc}}, W_{17}^{\mathrm{bc}},
W_{18}^{\mathrm{bc}}$ from~\eqref{eq:W-round-inv} at rounds
$14,15,16,17,18$ respectively (all from the backward chain and the context).

Finally, initialise $\hat{W}_9$ with $a_2 = a_3 = 0$:
\begin{equation}
  \hat{W}_9 = W_9^{\mathrm{base}} - T_1^{(5)}\big|_{a_3=0}
              - \Ch\!\left(e_8,\,T_1^{(7)}+0,\,T_1^{(6)}\big|_{a_2=a_3=0}\right),
  \label{eq:W9hat}
\end{equation}
where $T_1^{(5)}\big|_{a_3=0}$ denotes $T_1^{(5)}$ evaluated with
the sweep variable $a_3$ set to zero (and $a_2 = 0$).

\subsection{Phase C0: The Provisional-\texorpdfstring{$W_9$}{W9} Form}
\label{sec:C0}

The constraint C0 involves $W_9$, which depends on the unknowns
$a_1,a_2,a_3$.  The attack does not solve C0 exactly.  It solves C0 under a
provisional assignment for the part of $W_9$ that depends on $a_2,a_3$, and
tests that assignment afterwards.  The following proposition makes the
assignment precise.

\begin{proposition}[Provisional-$W_9$ form]
\label{lem:w9-diff}
Let $W_9$ be the message word determined by the state through
equation~\eqref{eq:W-round-inv} at $r=9$, and let $\hat{W}_9$ be the
quantity~\eqref{eq:W9hat}, computed from the context with $a_2=a_3=0$.
Then
\[
  \hat{W}_9 - W_9 \;=\; a_1 + \varepsilon(a_2,a_3) \pmod{2^{32}},
\]
where
\[
  \varepsilon(a_2,a_3) = \bigl[T_1^{(5)} - T_1^{(5)}\!\bigr|_{a_2=a_3=0}\bigr]
    + \bigl[\Ch(e_8,e_7,e_6) - \Ch\!\bigl(e_8,T_1^{(7)},T_1^{(6)}\!\bigr|_{a_2=a_3=0}\bigr)\bigr]
\]
depends on $a_2,a_3$ only through $\Maj(a_4,a_3,a_2)$, $\Maj(a_5,a_4,a_3)$
and the $\Ch$ arguments $e_6=a_2+T_1^{(6)}$ and $e_7=a_3+T_1^{(7)}$, and
vanishes at $a_2=a_3=0$.
\end{proposition}

\begin{proof}
Apply~\eqref{eq:W-round-inv} at $r=9$:
\[
  W_9 = \underbrace{T_1^{(9)} - \Sig{1}(e_8) - K_9}_{W_9^{\mathrm{base}}}
        - e_5 - \Ch(e_8,e_7,e_6).
\]
$T_1^{(9)}$ and $e_8$ depend only on the context $a_4,\ldots,a_9$, so
$W_9^{\mathrm{base}}$ is a context constant.  The identity
$e_r = a_{r-4} + T_1^{(r)}$ at $r = 5, 6, 7$ gives
\begin{align}
  e_5 &= a_1 + T_1^{(5)}, \label{eq:e5} \\
  e_6 &= a_2 + T_1^{(6)}, \label{eq:e6} \\
  e_7 &= a_3 + T_1^{(7)}, \label{eq:e7}
\end{align}
where $T_1^{(7)} = a_7 - T_2^{(7)}$ is a context constant
($T_2^{(7)} = \Sig{0}(a_6)+\Maj(a_6,a_5,a_4)$), $T_1^{(6)}$ depends on
$a_3$ through $\Maj(a_5,a_4,a_3)$, and $T_1^{(5)}$ depends on $a_2,a_3$
through $\Maj(a_4,a_3,a_2)$.  Substituting~\eqref{eq:e5},
\[
  W_9 = W_9^{\mathrm{base}} - a_1 - T_1^{(5)} - \Ch(e_8,e_7,e_6),
\]
while by definition~\eqref{eq:W9hat}
\[
  \hat{W}_9 = W_9^{\mathrm{base}} - T_1^{(5)}\!\bigr|_{a_2=a_3=0}
              - \Ch\!\bigl(e_8,\; T_1^{(7)},\; T_1^{(6)}\!\bigr|_{a_2=a_3=0}\bigr).
\]
Subtracting gives the stated expression.  The two brackets are the
$a_2,a_3$-dependent parts of $T_1^{(5)}$ and of the $\Ch$ term, and both
are zero at $a_2=a_3=0$.
\end{proof}

The residual $\varepsilon$ is not zero in general: on a random state it is
a nonzero 32-bit word.  The attack proceeds \emph{as if} $\varepsilon=0$.
We call this the \emph{provisional assignment}.  Under it, C0 reduces to a
single table lookup (Proposition~\ref{prop:c0}); C1 and C2 are then solved
for $a_2,a_3$ by iteration (\S\ref{sec:C1C2}); and finally
$\varepsilon(a_2,a_3)=0$ is tested exactly (\S\ref{sec:w9filter}).  A
candidate that passes the test satisfies all three constraints with the
true $W_9$, and forward evaluation confirms it.  The attack therefore finds
those preimages whose $(a_2,a_3)$ happen to satisfy $\varepsilon=0$.  Per context this is a $2^{-32}$ fraction of the roughly $2^{32}$ preimages that the three constraints leave over $(a_0,\ldots,a_3)$, i.e.\ about one preimage per context.  Whether the sweep reaches it is a separate question: the iteration is seeded on partial consistency with $\hat W_9$ and converges preferentially to $\varepsilon=0$ points, far above the $2^{-32}$ per fixed point that a uniform residual would give; how often is measured in \S\ref{sec:complexity}.  For
each of the seven verified preimages $\hat W_9 - W_9 = a_1$ holds exactly, as
it must for any solution the attack returns; for P1,
$\mathtt{0x72c05667} - \mathtt{0x6d1e5a20} = \mathtt{0x05a1fc47} = a_1$.
\paragraph{C0 lookup.}
For each $a_0 \in [0, 2^{32})$, compute:
\begin{align}
  e_0 &= a_0 + C_{e_0}, \label{eq:e0} \\
  W_0 &= a_0 + C_0, \label{eq:W0} \\
  g(a_0) &= -\Sig{0}(a_0) - \Maj(a_0,a_{-1},a_{-2})
            - e_{-3} - \Sig{1}(e_0) - \Ch(e_0,e_{-1},e_{-2}) - K_1,
  \label{eq:g} \\
  F_0 &= W_{16}^{\mathrm{bc}} - \sig{1}(W_{14}) - g(a_0)
         - \hat{W}_9 - W_0. \label{eq:F0}
\end{align}
Note that $g(a_0) = W_1 - a_1$: it is the portion of $W_1$ that
depends only on $a_0$ (derived from the round-1 inverse equation,
with $a_1$ factored out).  Then:

\begin{proposition}[C0 cancellation]
\label{prop:c0}
Under the provisional assignment $\varepsilon(a_2,a_3)=0$,
$F_0 = \sig{0}(W_1) - W_1 \pmod{2^{32}}$.
\end{proposition}
\begin{proof}
Under the provisional assignment, Proposition~\ref{lem:w9-diff} gives
$W_9 = \hat{W}_9 - a_1$, so from the C0 schedule constraint
\[
  \sig{0}(W_1)
  = W_{16}^{\mathrm{bc}} - \sig{1}(W_{14}) - W_9 - W_0
  = W_{16}^{\mathrm{bc}} - \sig{1}(W_{14}) - \hat{W}_9 + a_1 - W_0.
\]
Since $W_1 = a_1 + g(a_0)$:
\[
  \sig{0}(W_1) - W_1
  = W_{16}^{\mathrm{bc}} - \sig{1}(W_{14}) - \hat{W}_9 + a_1 - W_0
    - a_1 - g(a_0) = F_0.
\qedhere
\]
\end{proof}

The $a_1$ terms cancel exactly.  A single table lookup
$\mathcal{T}[F_0]$ returns $W_1$, and then $a_1 = W_1 - g(a_0)$.

\subsection{Phases C1/C2: Fixed-Point Iteration}
\label{sec:C1C2}

For each C0 survivor $(a_0, a_1, W_1)$, derive:
\begin{align}
  e_1 &= a_{-3} + a_1 - \Sig{0}(a_0) - \Maj(a_0,a_{-1},a_{-2}), \\
  T_2^{(2)} &= \Sig{0}(a_1) + \Maj(a_1,a_0,a_{-1}), \\
  F_{12} &= -T_2^{(2)} - e_{-2} - \Sig{1}(e_1) - \Ch(e_1,e_0,e_{-1}) - K_2.
\end{align}
Note that $F_{12} = W_2 - a_2$: the portion of $W_2$ depending only on
$(a_0, a_1)$.

\begin{proposition}[C1 and C2 cancellations]
\label{prop:c1c2}
Define the $a_3$-dependent quantity:
\[
  D(a_3) = \bigl(a_6 - \Sig{0}(a_5) - \Maj(a_5,a_4,a_3)\bigr)
            + \Ch\!\bigl(e_9,\, e_8,\, a_3 + T_1^{(7)}\bigr).
\]
Let $W_{16}^{\mathrm{sched}} = \sig{1}(W_{14}) + \hat{W}_9 + \sig{0}(W_1) + W_0$
and $W_{16}^* = W_{16}^{\mathrm{sched}} - a_1$.  Then:
\begin{align}
  F_{\mathrm{C1}}(a_3) &= W_{17}^{\mathrm{bc}} - \sig{1}(W_{15})
    - W_1 - \mathit{CONST}_{10} - F_{12} + D(a_3)
    = \sig{0}(W_2) - W_2, \tag{C1-target}\\
  F_{\mathrm{C2}}(a_2,a_3) &= W_{18}^{\mathrm{bc}} - \sig{1}(W_{16}^*)
    - W_{11}^{\mathrm{base}} - W_2 - F_{23}(a_0,a_1,a_2)
    = \sig{0}(W_3) - W_3, \tag{C2-target}
\end{align}
where $F_{23}(a_0,a_1,a_2)$ is the portion of $W_3$ depending on
$(a_0,a_1,a_2)$ but not $a_3$.  Both identities are unconditional for every C0 survivor: the lookup $\mathcal{T}[F_0]$ enforces $\sig{0}(W_1)-W_1=F_0$, and substituting the definition of $F_0$ gives $W_{16}^* = W_{16}^{\mathrm{bc}}$ exactly, so $\sig{1}(W_{16}^*)$ may equally be written $\sig{1}(W_{16}^{\mathrm{bc}})$.
\end{proposition}

\begin{proof}
\textbf{C1.}  From~\eqref{eq:W-round-inv} at $r=10$,
using $\mathit{CONST}_{10} = T_1^{(10)} - \Sig{1}(e_9) - K_{10}$:
\[
  W_{10} = \mathit{CONST}_{10} - e_6 - \Ch(e_9,e_8,e_7).
\]
By the shift-register identity $e_r = a_{r-4} + T_1^{(r)}$:
$e_6 = a_2 + T_1^{(6)}(a_3)$ and $e_7 = a_3 + T_1^{(7)}$, so
$\Ch(e_9,e_8,e_7) = \Ch(e_9,e_8,\,a_3+T_1^{(7)})$.
Grouping the $a_3$-dependent terms into $D(a_3)$:
\[
  W_{10} = \mathit{CONST}_{10} - a_2 - D(a_3).
\]
Substituting into the C1 schedule constraint
$\sig{0}(W_2) = W_{17}^{\mathrm{bc}} - \sig{1}(W_{15}) - W_{10} - W_1$:
\begin{align*}
  \sig{0}(W_2)
  &= W_{17}^{\mathrm{bc}} - \sig{1}(W_{15})
    - (\mathit{CONST}_{10} - a_2 - D(a_3)) - W_1 \\
  &= W_{17}^{\mathrm{bc}} - \sig{1}(W_{15}) - \mathit{CONST}_{10}
    + a_2 + D(a_3) - W_1.
\end{align*}
Using $W_2 = a_2 + F_{12}$, the $a_2$ terms cancel:
\begin{align*}
  \sig{0}(W_2) - W_2
  &= \bigl(W_{17}^{\mathrm{bc}} - \sig{1}(W_{15}) - \mathit{CONST}_{10}
     + a_2 + D(a_3) - W_1\bigr) - (a_2 + F_{12})\\
  &= W_{17}^{\mathrm{bc}} - \sig{1}(W_{15}) - W_1
     - \mathit{CONST}_{10} - F_{12} + D(a_3)
   = F_{\mathrm{C1}}(a_3). \qed_{\text{C1}}
\end{align*}

\textbf{C2.}  From~\eqref{eq:W-round-inv} at $r=11$:
$W_{11} = T_1^{(11)} - \Sig{1}(e_{10}) - \Ch(e_{10},e_9,e_8) - e_7 - K_{11}$.
Since $e_7 = a_3 + T_1^{(7)}$:
\[
  W_{11} = \underbrace{\bigl[T_1^{(11)} - T_1^{(7)} - \Sig{1}(e_{10})
             - \Ch(e_{10},e_9,e_8) - K_{11}\bigr]}_{W_{11}^{\mathrm{base}}} - a_3.
\]
So $W_{11} = W_{11}^{\mathrm{base}} - a_3$ with $W_{11}^{\mathrm{base}}$ determined
by context alone.

For a C0 survivor, $F_0 = \sig{0}(W_1) - W_1$, and unwinding the definition of $F_0$ gives $W_{16}^{\mathrm{bc}} = \sig{1}(W_{14}) + \hat W_9 + \sig{0}(W_1) + W_0 - a_1 = W_{16}^*$, so $\sig{1}(W_{16}) = \sig{1}(W_{16}^*)$ with $W_{16}=W_{16}^{\mathrm{bc}}$ the true round-16 word.
Substituting into C2:
$\sig{0}(W_3) = W_{18}^{\mathrm{bc}} - \sig{1}(W_{16}^*) - W_{11} - W_2$
$= W_{18}^{\mathrm{bc}} - \sig{1}(W_{16}^*) - (W_{11}^{\mathrm{base}} - a_3) - W_2$.
Using $W_3 = a_3 + F_{23}$, the $a_3$ terms cancel:
\begin{align*}
  \sig{0}(W_3) - W_3
  &= W_{18}^{\mathrm{bc}} - \sig{1}(W_{16}^*) - W_{11}^{\mathrm{base}}
     + a_3 - W_2 - a_3 - F_{23}\\
  &= W_{18}^{\mathrm{bc}} - \sig{1}(W_{16}^*) - W_{11}^{\mathrm{base}}
     - W_2 - F_{23}
   = F_{\mathrm{C2}}(a_2,a_3).
\end{align*}
\end{proof}

\paragraph{Iteration.}
The joint C1/C2 system defines a map
\[
  (a_2,a_3) \;\mapsto\;
  \bigl(\mathcal{T}[F_{\mathrm{C1}}(a_3)] - F_{12},\;
        \mathcal{T}[F_{\mathrm{C2}}(a_2,a_3)] - F_{23}\bigr).
\]
We iterate from a seed $(a_2^{(0)}, a_3^{(0)})$ for up to $8$ steps.
A \emph{convergent pair} is one at which both C1 and C2 are satisfied
simultaneously, i.e., the map is a fixed point.  Empirical convergence
rate: $1.3$--$1.9\times10^{-5}$ per C0 survivor (E1, \S\ref{sec:complexity}).

\paragraph{K-seed optimisation.}
Pre-computing $K$ initial pairs $(a_2,a_3)$ that satisfy the low-16-bit constraint $W_9^{\mathrm{conv}}(a_2,a_3) \bmod 2^{16} = \hat{W}_9 \bmod 2^{16}$, i.e.\ $\varepsilon(a_2,a_3) \equiv 0 \pmod{2^{16}}$, was measured to improve the convergence rate by $\approx 3.3\times$ on fixed points in an early pilot; no per-context lo-pass difference is visible between the $K_{\mathrm{seeds}}{=}0$ and $K_{\mathrm{seeds}}{=}4$ datasets of Table~\ref{tab:stats}.

\subsection{W9 Bit-Filter}
\label{sec:w9filter}

After each convergent pair $(a_2, a_3)$ is found, recompute
$\hat W_9 - \varepsilon(a_2,a_3)$, the $a_1$-free part of $W_9$ at the
converged values:
\begin{equation}
  W_9^{\mathrm{conv}} = W_9^{\mathrm{base}}
    - \bigl(a_5 - \Sig{0}(a_4) - \Maj(a_4,a_3,a_2)\bigr)
    - \Ch\!\bigl(e_8,\, a_3+T_1^{(7)},\, a_2+T_1^{(6)}(a_3)\bigr).
  \label{eq:W9true}
\end{equation}

Two sequential filters are applied:
\begin{itemize}
  \item \textbf{Lo filter:}
    $(W_9^{\mathrm{conv}} \bmod 2^{16}) = (\hat{W}_9 \bmod 2^{16})$,
    i.e.\ the low half of $\varepsilon$ vanishes.
    Passes at $52$--$85\times$ (E1) or $37$--$44\times$ (E2) the uniform
    $2^{-16}$ rate (experiments defined in \S\ref{sec:complexity}).
  \item \textbf{Hi filter:}
    $(W_9^{\mathrm{conv}} \gg 16) = (\hat{W}_9 \gg 16)$, i.e.\ the high
    half vanishes too, so that both together test $\varepsilon=0$ exactly.
    Passes at $12.6$--$13.3\times$ (E1) or $7$--$11\times$ (E2) the uniform
    rate.  This filter is
    \emph{not} uniform; a coarse-binned test suggests otherwise but is
    underpowered against a spike at a single residual (\S\ref{sec:chi2}).
\end{itemize}
The combined rate is not obtained by multiplying the two marginals: the
iteration is seeded on lo-consistency, so independence of the two halves
cannot be assumed and has not been established (\S\ref{sec:complexity}).  Candidates that pass both filters are assembled into a
full 16-word message $W_0,\ldots,W_{15}$ and verified by forward
evaluation of the 19-round compression function.

\subsection{Complexity Analysis}
\label{sec:complexity}

\begin{table}[ht]
\centering
\caption{Empirical per-context statistics.  E1: two deduplicated CPU
pilots of $0.20$ and $1.37$ context-equivalents.  Lo-pass counts per
context: the E1 pilots, the 165-context H100 run of \S\ref{sec:chi2}, and
the campaign control arm of \S\ref{sec:screening}.  $\mu$ rows: the
240-context campaign.  Wall times: H100.}
\label{tab:stats}
\small
\begin{tabular}{@{}lr@{}}
\toprule
Quantity & Value \\
\midrule
C0 survivors (table hits) & $2{,}721{,}603{,}628$ ($63.37\%$ of $2^{32}$) \\
\multicolumn{2}{@{}l}{\emph{\textbf{E1} --- deduplicated (unique fixed points);
two CPU pilots}}\\
C1/C2 fixed points & $3.6$--$5.1\times10^{4}$ per context ($1.3$--$1.9\times10^{-5}$ per C0 survivor) \\
$W_9$ lo-pass events per context & $28$ (E1), $35.0$ (165-context run), $33.7$ (campaign control) \\
$W_9$ lo-pass rate per fixed point & $7.9\times10^{-4}$--$1.3\times10^{-3}$ \\
Structural \emph{lo} enhancement & $52$--$85\times$ above $2^{-16}$ \\
$W_9$ hi-pass rate per fixed point & $\approx 2.0\times10^{-4}$ \\
Structural \emph{hi} enhancement & $12.6$--$13.3\times$ above $2^{-16}$ \\
\midrule
\multicolumn{2}{@{}l}{\emph{Legacy trajectory counts (each fixed point counted
once per remaining iteration)}}\\
Repeat-count inflation factor & $5.5$--$5.9\times$ (pilots), $6.35\times$ (campaign) \\
\midrule
$\mu$ (original model: measured lo count $\times$ & \\
\quad uniform conditional hi) & $\approx 5\times 10^{-4}$ \\
$\mu$ (directly measured, 240 ctx) & $1.25\times10^{-2}$\ \ CI $[2.6{\cdot}10^{-3},3.7{\cdot}10^{-2}]$ \\
\quad screened arm (120 ctx, $N_{LH}=3$) & $2.50\times10^{-2}$ \\
\quad control arm (120 ctx, $N_{LH}=0$) & $<3.1\times10^{-2}$ (exact Poisson, 95\%) \\
$\mu$ (independence-derived projection) & $6\times10^{-3}$--$1.3\times10^{-2}$ \\
Observed contexts (P3, P2, P1) & 4,\ 21,\ 35\ \ (mean 20) \\
Wall time per context (H100), $K_{\mathrm{seeds}}{=}0$ & $\approx 25\;\mathrm{s}$ \\
\quad same, $K_{\mathrm{seeds}}{=}4$ (campaign) & $\approx 89\;\mathrm{s}$ \\
\bottomrule
\end{tabular}
\end{table}

\paragraph{W9 structural enhancement.}
Neither half of the $W_9$ filter is uniform.  Both enhancements arise because
the C1/C2 fixed-point iteration, started from $W_9$-lo-consistent seed pairs,
preferentially converges to fixed points that also satisfy the $W_9$
constraint---an effect that acts on the full 32-bit residual, not only on the
low half.  Experiment \textbf{E2}, an independent CPU reimplementation
measuring over $2.2\times10^{9}$ swept $a_0$ values per configuration, gives:

\begin{itemize}\itemsep2pt
\item the \emph{lo} filter passes at $37$--$44\times$ the uniform $2^{-16}$
      rate, confirming the enhancement mechanism;
\item the \emph{hi} filter passes at $7$--$11\times$ the uniform rate---it
      is \emph{not} uniform, contrary to what a coarse-binned test suggests
      (\S\ref{sec:chi2}).
\end{itemize}

E1 and E2 are separate campaigns and their rates are not directly comparable:
E1 deduplicates fixed points and reports per-fixed-point rates from two CPU
pilots, whereas E2 reports per-swept-$a_0$ rates from a distinct
implementation on different contexts.  Both find the same qualitative effect
--- a large \emph{lo} enhancement and a smaller but nonzero \emph{hi}
enhancement --- and the residual factor of roughly $1.5$ between them is not
currently explained; it is within the spread expected from context-to-context
variation (\S\ref{sec:screening}), which is itself large.  \textbf{E1 is the
headline dataset}, since it is deduplicated and feeds the $\mu$ projection.
Neither E1 nor E2 measures the success rate; the preregistered campaign
below does.

These rates are not artefacts of the lookup table's representative policy.
A third experiment (\textbf{E3}) repeated the measurement on byte-identical
contexts with a table built to keep the \emph{smallest} rather than the
\emph{largest} preimage of each target: lo enhancement $30.0\times$
(smallest) against $28.5\times$ (largest), and hi $11.1\times$ against
$6.9\times$, agreeing within Poisson error on the $8$ and $13$ hi events
observed.  E3 used its own contexts and sample, which is why its lo figure
sits below E1's range.  Direct tests of lo/hi independence at reduced filter widths give
$D = P(\text{both})/[P(\text{lo})P(\text{hi})] = 1.00\pm0.03$ at $4$ bits and
$1.06\pm0.10$ at $6$ bits, consistent with independence at the widths where
joint events can be counted.

Multiplying the measured marginals by the deduplicated fixed-point count
gives a \emph{projection} of $\mu \approx 6\times10^{-3}$--$1.3\times10^{-2}$
per context (one figure per pilot).  Once fixed points are deduplicated the
lo-pass count is $28$--$35$ per context across the pilots, the 165-context
run and the campaign control arm (Table~\ref{tab:stats}), so the lo
analysis stands; the correction is almost entirely the hi filter, worth about
$13\times$ rather than the two orders of magnitude a naive comparison of
non-deduplicated counts suggests.  We
label this a projection rather than a measurement for two reasons.  First, it
assumes the two halves are independent; since the iteration is seeded on
lo-consistency, that is exactly the assumption one should not make for free.
Half-width diagnostics detect no material dependence at 4 and 6 bits, but a
sparse spike at the exact 16-bit value would remain invisible to them---the
same limitation that makes the binned $\chi^2$ test uninformative
(\S\ref{sec:chi2}).  Second, the direct estimator $\hat\mu = N_{LH}/C$ now has
events to work with, though few.  The preregistered campaign of
\S\ref{sec:screening} attacked 240 contexts at full $2^{32}$ exposure and
recorded $N_{LH}=3$ complete 32-bit matches, all of which verified as
preimages, giving
\[
  \hat\mu \;=\; \frac{3}{240} \;=\; 1.25\times10^{-2}\ \text{per context},
  \qquad 95\%\ \text{CI}\ [2.58\times10^{-3},\ 3.65\times10^{-2}].
\]
Split by arm, the screened half gave $3/120 = 2.50\times10^{-2}$ and the
control half $0/120$, i.e.\ $\mu < 3.1\times10^{-2}$ (upper limit of the
exact two-sided 95\% Poisson interval for zero events in 120 contexts).  Two cautions apply.  The pooled figure mixes the two arms and the
screened half is enriched by construction, so it is not an estimate for
uniformly sampled contexts; the control arm is the closer analogue and it
produced no events.  And three events give an interval spanning a factor of
$14$, so this bounds the rate rather than pinning it.  What it does settle is
that the rate is not pathologically below the projection: the measured
$1.25\times10^{-2}$ falls inside the $6\times10^{-3}$--$1.3\times10^{-2}$ band
projected before the campaign was run, which is a genuine out-of-sample
check on the independence-derived model rather than a fit to it.
The projections themselves remain unstable.  After adopting
global fixed-point deduplication, the independence-derived projection was
$1.31\times10^{-2}$ in a $0.20$ context-equivalent pilot and
$5.80\times10^{-3}$ in a $1.37$ context-equivalent pilot.  That variation
shows the marginal-event sample is too small for a stable projection; it does
not by itself establish a directional bias.  An earlier figure of
$4.88\times10^{-2}$ used repeat-counted fixed points and is not comparable
with either.

Two earlier sources of error are worth recording, since both inflate
$\mu$ if left uncorrected.  A converged lane remains at its fixed point for
every subsequent iteration, so counting convergences per iteration overcounts
each fixed point---measured at $5.5$--$5.9\times$ in the CPU pilots and $6.35\times$ over the 240 campaign contexts (the \texttt{fp\_replay} and \texttt{fp\_trajectories} columns of the campaign data).  And
convergences tallied per seed chain rather than per context understate the
per-context count by the number of chains.  All figures above are globally
deduplicated on the fixed-point identity $(a_0,a_1,a_2,a_3)$ within a context.

\paragraph{What the three dedicated finds do and do not establish.}
The three dedicated runs found preimages in their 4th, 21st and 35th contexts (solver indices 3, 20 and 34): three successes over 60 attempted contexts, $\hat p = 0.050$ with an exact 95\% interval of $[0.010,\,0.139]$.  The $6\times10^{-3}$ projection lies \emph{below} that interval. Under $p=0.006$, observing three or more successes in 60 contexts has probability $0.006$, so either the historical runs were unusually favourable
or the projection is low.  Two very small samples are in tension here, and
neither settles the rate.  It is not compatible with the uniform-filter model: at
$p\approx5\times10^{-4}$ the same observation has probability $4\times10^{-6}$,
so the finds do falsify the uniform model even though they cannot pin the true
rate.  We therefore present them as demonstrations of feasibility, not as a
runtime estimator.  Three observations cannot support an expected-value claim,
and the per-context success probability $\Pr(N_{LH}\ge1)$---which governs time
to first preimage---is not in general $\mu$ itself if successful fixed points
cluster in favourable contexts.  The preregistered multi-target campaign of \S\ref{sec:screening} reports both:
over 240 fully attacked contexts, $\hat\mu = 3/240 = 1.25\times10^{-2}$ and
$\widehat{\Pr}(N_{LH}\ge1) = 3/240$ also, no context having produced more than
one match.  With three events these coincide and neither is tightly
determined.  The combined search time for the
three was $1{,}464$\,s on the H100 (871.6, 511.1 and 81.7\,s), i.e.\
${\approx}25.7$\,s per full sweep (57 complete sweeps preceded the three hits) with the single $(0,0)$ start.  We note for
comparison that the screening campaign (\S\ref{sec:screening}), which uses
$K_{\mathrm{seeds}}{=}4$ and therefore runs four seed chains per context,
measured $89$\,s per context over 240 contexts.  The ratio $3.5$ is consistent
with the $4\times$ increase in inner-loop work and not with any difference in
the sweep itself; per-context time was flat in the fixed-point count
($90.2$\,s in the screened arm against $88.3$\,s in the control arm despite
$8.1\times$ the fixed points), confirming the cost is the $2^{32}$ sweep rather
than the downstream bookkeeping.

\subsection{Context Screening}
\label{sec:screening}

The clustering just hypothesised is real and is large enough to exploit.  The
attack samples a context and then sweeps $a_0$ over $2^{32}$, giving every
context the same budget; that is optimal only if contexts are interchangeable.
Over 280 contexts, each scored on a $2^{21}$ sample of $a_0$, the deduplicated C1/C2 fixed-point count had mean $63.4$ and variance $97{,}082$, a variance-to-mean ratio of
${\approx}1{,}530$ against the $1.0$ that a Poisson model would give, with
median $12$ and maximum $4{,}924$.  The top decile of contexts held $68.6\%$ of
the total yield.

Overdispersion alone is not usable, since it could be per-sample noise.  The
operative question is whether a cheap score measured on one $a_0$ sample
predicts yield on an independent one.  Scoring each of 392 contexts twice on
independent $2^{20}$ samples gave a Spearman rank correlation of $0.929$ in the archived run (\texttt{screening\_validation.py}, seed 99001, \texttt{data\_screening\_validation.json}; $0.910$ in an earlier unarchived run) between the two scores, so productivity is a stable property of the context and
is recoverable from a short pass.  Since screening costs $2^{20}$ against a full
sweep of $2^{32}$, it is $1/4096$ of a context's cost and is nearly free:
selecting the top $1\%$ on the screening sample gave a $12.99\times$
out-of-sample \emph{fixed-point} yield lift for a $2.44\%$ overhead.  We stress
that this is a statement about fixed points only; the attack-relevant lift is
measured below and is substantially smaller.  The deep tail is also unstable
---the top-$1\%$ row rests on 4 contexts of 392 and moved from $12.99\times$ to
$31.93\times$ between runs---so the top-decile figure ($4.98\times$ and $6.79\times$ in two runs, of which only the second is archived) is the defensible one at this level.

\paragraph{Does the lift survive conversion?}
The score counts C1/C2 fixed points, so the lift above is in the first place a
fixed-point lift.  Whether it survives to the events that actually determine
attack cost is a separate question, and we settled it with a preregistered
paired campaign (\texttt{campaign\_screening.py}, seed 20260810; the run manifest in \texttt{data\_campaign\_screening.json} records the command, seeds and software versions; it records the code commit as ``unknown'', identifies the 40 targets by index and a truncated digest, and does not store the random messages that generated them, so the campaign is reproducible from the recorded seed and command but not from the manifest alone) rather than by assuming proportionality.

\emph{Design, fixed before the run.}  40 independent targets; 10 contexts
generated per target; every context scored on an identical cheap $2^{20}$ $a_0$
sample; the top $30\%$ by score forming the \textsc{screened} arm and an equal
number drawn at random from the remainder forming the \textsc{control} arm;
both arms then given \emph{exactly} the same full $2^{32}$ exposure per
context, with no early stop on a hit, so that a lucky context is never
under-exposed.  240 contexts were fully attacked, 120 per arm, in $5.95$\,h of attack time
on one H100 (screening excluded).

\begin{center}
\begin{tabular}{@{}lrrr@{}}
\toprule
Counter (per context) & screened & control & lift \\
\midrule
C1/C2 fixed points & $159{,}545$ & $19{,}624$ & $8.13\times$ \\
$W_9$ lo-passes    & $114.86$    & $33.73$    & $\mathbf{3.41\times}$ \\
$W_9$ hi-passes (given lo) & $0.025$ & $0$ & --- \\
full 32-bit $W_9$  & $0.025$     & $0$        & --- \\
verified preimages & $0.025$     & $0$        & --- \\
\bottomrule
\end{tabular}
\end{center}

\emph{The conversion is sublinear.}  An $8.13\times$ lift in fixed points buys
only a $3.41\times$ lift in lo-passes: screened contexts yield more fixed
points, but each individual fixed point is \emph{less} likely to pass the lo
filter.  An earlier pilot on 10 lo events suggested the conversion was
proportional; that pilot was underpowered and this campaign supersedes it.

\emph{The effect is real but highly heterogeneous.}  A target-blocked bootstrap
(20{,}000 resamples, seed 20260810, percentile interval, blocking on target because contexts within a target are not independent) gives a $95\%$ interval of $[2.13\times, 5.82\times]$ on the
pooled lo lift.  Per target the median lift is $3.78\times$ and the geometric
mean $3.80\times$, but the interquartile range (Weibull quantiles) spans
$1.31\times$ to $12.35\times$: screening helped on $31$ of $40$ targets and \emph{hurt} on the
other $9$, where it selected contexts richer in fixed points but poorer in
lo-passes.  A distribution-free sign test on the per-target lifts---which is
insensitive to the heavy tail that makes the pooled ratio unstable---gives $p=3.4\times10^{-4}$ one-sided ($6.8\times10^{-4}$ two-sided).

\emph{How the counter counts.}  The lo-pass counter is incremented once per converged trajectory: the four seed chains are deduplicated within a chain but not across chains, so a fixed point reached by two chains is counted twice.  Duplicated fixed points are $0.034\%$ of screened-arm trajectories ($6{,}541$ of $19{,}151{,}936$; $4{,}053$ of them in a single context) and $0.002\%$ of control-arm trajectories ($42$ of $2{,}354{,}867$); their lo status was not recorded.  Because the lo test is a deterministic function of the fixed point $(a_2,a_3)$, a duplicate has the same lo status as its original, so a context can lose at most $\min(\text{duplicates},\text{lo-passes})$ lo-passes under identity-level deduplication: $3{,}972$ in the screened arm and $42$ in control.  That bounds the deduplicated pooled lift within $[2.42\times, 3.44\times]$, with point estimate $3.405\times$ if duplicated points pass the lo filter at the ordinary rate; since seed chains are constructed to be $W_9$-consistent this is not guaranteed, and the campaign was not rerun with cross-chain deduplication.  In the artefact the screening scorer (\texttt{context\_screening.py}) now deduplicates across chains; the solver that produced the campaign counts (\texttt{h100\_extended.py}) is left as it ran, so the published counts can be regenerated exactly, and it records the unique fixed-point count separately (\texttt{fixed\_points} against \texttt{fp\_trajectories} per row).

\emph{What remains unmeasured.}  The screened arm produced 3 full 32-bit $W_9$
matches and 3 verified preimages against 0 in control.  Under the null of
equal rates each event falls in the screened arm with probability $1/2$,
so all three doing so has $p = 1/8 = 0.125$ (one-sided): suggestive and
\emph{not} significant.  We therefore
report no lift for the full-match or preimage counters.  Establishing one would
require of order $10^2$ GPU-hours, which we did not spend.

\emph{Summary of the claim.}  Screening costs $1/4096$ of a context's budget
and delivers a measured $3.41\times$ ($95\%$ CI $[2.13, 5.82]$) improvement in
lo-pass yield per attacked context.  It is \textbf{not} the $12.99\times$
figure quoted for fixed points, which we retain above only as a statement about
fixed points; and it is not yet established as a speedup in time-to-preimage.

%-----------------------------------------------------------------------
\section{Experimental Results}
\label{sec:experiments}
%-----------------------------------------------------------------------

\subsection{Hardware and Software}

All GPU experiments were run on NVIDIA H100 SXM5 80\,GB HBM3 via RunPod
cloud compute; the E1 and E2 pilots ran on CPU.  The attack code uses
PyTorch~2.x on CUDA, with the C0 sweep vectorised in batches of $2^{24}$
elements.  The 16\,GB $\sigma_0$-difference table is held in GPU memory
throughout the
sweep; table lookups are implemented as indexed tensor reads.

\subsection{Verified Preimages}

\begin{table}[p]
\centering
\caption{Verified 19-round preimages.  P1--P3 were found in dedicated runs
($K_{\mathrm{seeds}}{=}0$); P4--P6 in the screened arm of the campaign of
\S\ref{sec:screening} ($K_{\mathrm{seeds}}{=}4$; target and context number
given).  P7 is a preimage of the all-ones digest, the target of
Zaikin~\cite{Zaikin2026} (\S\ref{sec:comparison}), found in a dedicated run
with $K_{\mathrm{seeds}}{=}4$.  Only the target digest was supplied to the solver; the six random targets are 19-round digests of random 55-byte messages, so a preimage was known to exist, whereas the all-ones digest is the structured target of~\cite{Zaikin2026}, not known in advance to have one.  All values are hexadecimal.  Context numbers are the solver's zero-based indices.  $^\dagger$Found 6\,s into the fourth context (index 3).  $^\ddagger$Found 20\,s into the 84th context; 7{,}373\,s is the
whole run.}
\label{tab:preimages}
\footnotesize
\begin{tabular}{@{}lll@{}}
\toprule
 & \textbf{Hash / words} & \textbf{Time / note} \\
\midrule
P1 & \texttt{1e65261c 54255188 604f5375 09183973} & 871.6\,s \\
   & \texttt{3de63e96 6b5e4715 658226bf 03588447} & ctx~34 \\
   & \texttt{22f091af ec52d67b 74c33819 a280dc6a} & $W_0$--$W_7$ \\
   & \texttt{b001ff1a 1f2356a5 3eccf108 bd9a2333} & \\
   & \texttt{abe611d1 6d1e5a20 8041df25 e43d31af} & $W_8$--$W_{15}$ \\
   & \texttt{aa895a2e 69106ad2 7479fa3a 2a9abb91} & \\
\midrule
P2 & \texttt{fb52f81b aed24f87 28faf5bb ce82c67d} & 511.1\,s \\
   & \texttt{51076117 2fb9876d 9e3a72dd a351b7ca} & ctx~20 \\
   & \texttt{37e6702f bc20efea 2dd42a3e 501dfbe9} & $W_0$--$W_7$ \\
   & \texttt{3cacc578 ea2de1c1 11c0f066 0f22be47} & \\
   & \texttt{2a447d2d 13f0080f 1f33df6b d655d8e6} & $W_8$--$W_{15}$ \\
   & \texttt{15730eaa 9bf64950 9f129973 5a964edf} & \\
\midrule
P3 & \texttt{1bd7ebbd c4d938fb 26d19b5d d5caf333} & 81.7\,s \\
   & \texttt{de397bd1 c745727b d5556baf 38ccf977} & ctx~3$^\dagger$ \\
   & \texttt{3ce8fba4 e2fb9661 44730c59 e1cf4bc0} & $W_0$--$W_7$ \\
   & \texttt{e1a18d93 97658983 67efe2a7 ef260ecb} & \\
   & \texttt{d4c6dbe0 13e9388e 95664a59 4d9e248b} & $W_8$--$W_{15}$ \\
   & \texttt{74137862 664815ac 89eae95a cd7dbef5} & \\
\midrule
P4 & \texttt{c55aecaf d68ac4f5 59bb2b94 0688f3bc} & campaign \\
   & \texttt{1dabe3ff 86a7d034 47e8b11a 637d7506} & target 2, ctx~9 \\
   & \texttt{3694ee67 a1d020b9 47f08ea3 dc873e2b} & $W_0$--$W_7$ \\
   & \texttt{1f41c8b1 39018970 d3a2afb2 23f6e6cf} & \\
   & \texttt{80c46545 f7b2b46c 2e44f333 2effe68b} & $W_8$--$W_{15}$ \\
   & \texttt{06d80cdf 8525e7d7 7ef008e5 13e56087} & \\
\midrule
P5 & \texttt{a8241289 b5980304 e9dfd401 b1370200} & campaign \\
   & \texttt{ab32b5e8 cfaff26c 26aabcfc 8be89df4} & target 20, ctx~7 \\
   & \texttt{3e0a7d67 df0944a3 4bfe156b bd344ac5} & $W_0$--$W_7$ \\
   & \texttt{7b7d9587 8984e286 21c995d9 d8e53cdd} & \\
   & \texttt{0a0be79d e02dda1a e17e0f4a e3ff692b} & $W_8$--$W_{15}$ \\
   & \texttt{8c782631 97dc4811 5ed68656 dfce314f} & \\
\midrule
P6 & \texttt{09156002 031fed28 99bc67d0 e419b6ff} & campaign \\
   & \texttt{81b9a0dd 4e3fa417 e000817a db85b143} & target 25, ctx~4 \\
   & \texttt{164084a5 f97bd8bd bff47431 a8e55b29} & $W_0$--$W_7$ \\
   & \texttt{948d1160 10eb8f7d d0ea565e 7918b8b9} & \\
   & \texttt{87527a3a 6baa5685 a1fbc5dc 7548adfc} & $W_8$--$W_{15}$ \\
   & \texttt{be3fb206 b9738756 b4ca33e3 b8354a95} & \\
\midrule
P7 & \texttt{ffffffff ffffffff ffffffff ffffffff} & 7{,}373\,s \\
   & \texttt{ffffffff ffffffff ffffffff ffffffff} & ctx~84$^\ddagger$ \\
   & \texttt{3c6ff45c a5b23f8d dc6b8089 4f232935} & $W_0$--$W_7$ \\
   & \texttt{79578e14 9c2da339 b22ec37a 9554191b} & \\
   & \texttt{e69661f0 49a43b7a d8c60a2b 88d7588c} & $W_8$--$W_{15}$ \\
   & \texttt{6225084d 87c95c4e 23ee573d 8d8fedfa} & \\
\bottomrule
\end{tabular}
\end{table}

Each row of Table~\ref{tab:preimages} was independently verified by
the following procedure:
\begin{enumerate}
  \item Parse the 16 words $W_0,\ldots,W_{15}$ as a 512-bit block.
  \item Run the standard SHA-256 message schedule for 19 rounds,
    initialised with the standard IV.
  \item Add the final state to the IV word-by-word.
  \item Confirm the resulting 32-byte value equals the target hash.
\end{enumerate}
All seven preimages pass this check.  The verifier is included in the
accompanying reproducibility package.

\subsection{Chi-Squared Uniformity Test}
\label{sec:chi2}

To test the $1/2^{16}$ model for the W9-hi filter, we collected
the values $(W_9^{\mathrm{conv}} \gg 16) \bmod 2^{16}$ for all
lo-pass events across 165 GPU contexts (a total of $5{,}775$ lo-pass
events).  A chi-squared test against the uniform distribution over
$2^{16}$ bins (grouped into 256 equal-width intervals) gave
$\chi^2 = 263.4$ with 255 degrees of freedom ($p = 0.35$).

This test is \emph{not} sensitive to the effect that matters, and we record the
point because it is easy to be misled by it.  The enhancement is concentrated
at and near the exact 16-bit residual, whereas grouping into 256
equal-width intervals spreads it across a bin of width 256.  We measured
the dilution directly on the \emph{lo} residual, where the exact-value
enhancement is known: a $37\times$ enhancement at the exact residual
appears as $3.9\times$ at 256-bin resolution, a dilution of $9.5\times$.
A hi enhancement of $13\times$ diluted by the same factor would appear as
$1.4\times$, shifting the occupancy of one bin from $22.6$ to roughly
$31$---invisible to a $\chi^2$ statistic over 255 degrees of freedom.

Measuring at exact resolution instead shows the hi filter passing at
$7$--$13\times$ the uniform rate (E1 and E2, \S\ref{sec:complexity}).  The correct
conclusion is that the coarse-binned test is underpowered, not that the hi
residuals are uniform; filter rates in this attack must be estimated at exact
residual resolution.

%-----------------------------------------------------------------------
\section{Comparison with Prior Work}
\label{sec:comparison}
%-----------------------------------------------------------------------

\begin{table}[htbp]
\centering
\caption{Preimage results on step-reduced SHA-256.  ``Concrete'' means at
least one preimage was computed and verified.}
\label{tab:comparison}
\footnotesize
\begin{tabular}{@{}lllll@{}}
\toprule
Reference & Rounds & Cost & Target & Concrete \\
\midrule
Isobe--Shibutani~\cite{IsobeShibutani2009} & 24 & $2^{240}$ & any & no \\
Aoki et al.~\cite{AokiGuo2009} & 43 & $2^{254.9}$ & any & no \\
Khovratovich et al.~\cite{KRS2012} & 45 & $2^{255.5}$ & any & no \\
Guo et al.~\cite{GuoLiLiuWangZhang2026} & 47 & $2^{255.6}$ & any & no \\
Davydov et al.~\cite{Davydov2025} & 17, 18 & SAT, seconds & zero digest & yes \\
Davydov et al.~\cite{Davydov2025} & 19 & SAT, unsolved$^{a}$ & zero digest & partial \\
Zaikin~\cite{Zaikin2026} & 19 & $18$\,h\,$33$\,min on 192 CPU cores$^{b}$ & all-ones digest & yes (one) \\
\midrule
\textbf{This work} & \textbf{19} & ${\approx}80$ attempts of $2^{32}$ (${\approx}2$\,h, one H100)$^{c}$ & random (six); all-ones & \textbf{seven} \\
\bottomrule
\end{tabular}

\smallskip
\begin{minipage}{0.92\linewidth}\footnotesize
$^{a}$\,Reached only in a weakened form that constrains at most 31 of the
256 output bits.\quad
$^{b}$\,One run.  Parameterized Kissat inside Cube-and-Conquer, via an intermediate problem between rounds 18 and 19 whose solutions are 19-round preimages with probability $1/8$ by his account, so the expected cost of that route is about eight times this figure (see text).\quad $^{c}$\,Pooled rate of the 240-attempt campaign of \S\ref{sec:screening}: screened arm $3/120$, control arm $0/120$; 95\% interval $[2.6\times10^{-3},3.7\times10^{-2}]$ per attempt.  The six random targets are digests of random messages, so a preimage was known to exist.
\end{minipage}
\end{table}

The theoretical attacks reach far more rounds than any practical method,
at costs within a few bits of the generic $2^{256}$, and have not been
executed.  The practical line is SAT-based.  Davydov et
al.~\cite{Davydov2025} use three encoding tools, CBMC, Transalg and SAT-encoding, the last in its Counter-chain, Dot-matrix, Two-operand and Espresso adder variants, six CNF encodings in all, with the Kissat solver at a $5000$\,s limit, find preimages of
the 17- and 18-round compression function (18 rounds in $2.38$--$20.6$\,s
depending on encoding), and report 19 rounds as not solved for every
encoding tried.  Their weakened 19-round problem constrains only a zero
prefix of the output: plain Kissat reached a 24-bit prefix, and a 24-hour
cube-and-conquer run on 16 cores solved every case except the 32-bit one,
i.e.\ at most 31 of the 256 output bits.  The authors conclude that finding a 19-round preimage ``is too difficult a task even if the SAT encoding is carefully chosen and parallel computing is used'' (our translation; the article is in Russian with an English abstract).  Their published benchmark
instances fix the all-zero digest as the target.

Zaikin~\cite{Zaikin2026} inverted the 19-round compression function first,
for the all-ones digest with all 256 output bits fixed.  His method
constructs weakened intermediate problems between consecutive round counts,
tunes Kissat on them with ParamILS, and runs a Cube-and-Conquer solver with
the tuned Kissat.  Direct attempts on the 19-round instance found nothing
in seven days on 192 cores; the preimage came from the intermediate problem between rounds 18 and 19 that keeps 29 of the 32 bits of the message word entering round 19 and fixes the other three to constants, whose solutions are full 19-round preimages with probability $1/8$,
after $18$\,h\,$33$\,min on 192 CPU cores, about $3{,}560$ core-hours (his
Tables~8 and~9), against $25$\,s and $1$\,s on one core for 18 and 17
rounds.  Eight off-the-shelf parallel SAT solvers, seven from the SAT Competition 2024 and ALIAS, found nothing.  Since,
by his account, a solution of that intermediate problem is a 19-round
preimage with probability $1/8$ and the solutions found at $20$, $24$ and
$28$ kept bits were not, the reported time is a single favourable draw;
taken as an expectation his route costs about eight times that figure, and
there is no second success to estimate the spread.  Our rate, by contrast, is measured over $240$ preregistered attempts with an interval (\S\ref{sec:complexity}), and the $98$ further attempts on his target (two full matches, one of them discarded by the bug described below) and the $60$ of the three dedicated runs are consistent with it.  The attack model is the same as ours (his published preimage, Table~7 of~\cite{Zaikin2026}, verifies under \texttt{verify\_r19.py} at 19 rounds).  What the present work adds is a
different method with a different cost profile: after a one-time $16$\,GB
precomputation, each attempt is one $2^{32}$ sweep on a single GPU, about $25$\,s with a single seed chain and $89$\,s with the four chains used in the campaign of \S\ref{sec:screening}, where about one attempt in eighty succeeded when both arms are pooled ($3/120$ screened, $0/120$ control), so a preimage of an arbitrary digest costs about two hours of one H100 at that pooled rate, within a 95\% interval spanning a factor of 14, and seven preimages, six of random digests and one of the all-ones digest, are published together with the code that found them.  We do not convert
between GPU-hours and CPU core-hours; the two figures are simply stated.

\paragraph{Head-to-head on the same target.}
After reading~\cite{Zaikin2026} we ran the attack on its target, the
all-ones digest, with the campaign configuration ($K_{\mathrm{seeds}}{=}4$)
on one H100.  A first run produced a full 32-bit $W_9$ match at its 14th
context that the candidate-assembly code then discarded by crashing: it
indexed the compacted list of full matches with a position from the longer
list of lo-passes, which is correct only when the winner is first in its
batch, as all six earlier preimages had happened to be.  The bug is fixed
in the artefacts and the published counts are unaffected, since a crash
would have aborted the campaign.  The rerun with the fixed code found a
preimage 20\,s into its 84th context, after $7{,}373$\,s ($122.9$\,min) of
solving; it is P7 in Table~\ref{tab:preimages}, verified like the others,
and it differs from the preimage in Zaikin's Table~7, so it is also a
second preimage of that digest.  Counting the discarded match, the target
gave two full matches in 98 contexts, in line with the campaign rate.  The
raw log is \texttt{data\_ones\_target\_run.txt}.

\paragraph{Attack scope.}
The attack covers 19 of 64 rounds and targets the compression function
with the standard chaining value; the padded hash is discussed in
\S\ref{sec:intro}.  Full SHA-256 is not threatened, and nothing here bears
on its security margin.

%-----------------------------------------------------------------------
\section{Conclusion}
\label{sec:conclusion}
%-----------------------------------------------------------------------

Three cancellation identities in the SHA-256 round and schedule equations,
together with a provisional assignment for the part of $W_9$ that depends
on $a_2,a_3$, reduce the recovery of a 19-round preimage to one $2^{32}$
sweep of table lookups per attempt.  Six preimages of random digests and
one of the all-ones digest attacked by Zaikin~\cite{Zaikin2026} were found
and verified, and a preregistered 240-attempt campaign measured the
success rate directly, $\hat\mu = 1.25\times10^{-2}$ per attempt with a
wide interval.  A cheap screening pass improves the intermediate yield by
a measured $3.41\times$; whether that survives to the final success rate
is not established.

The same construction does not reach 20 rounds at the 19-round cost.  The fourth schedule
constraint adds a 32-bit condition, and the companion
note~\cite{Basra2026R20} shows why no rearrangement of this attack family
absorbs it: the full-avalanche dependency graph among the absorbed state
words acquires its first cycle at exactly 20 rounds, through the term
$\Sigma_0(a_4)$ entering $W_9$, so the 32-bit iteration that closes the
19-round system has no cheap counterpart.  We do not claim that this is
fundamental.

\paragraph{Artefacts.}
Solver, table construction, campaign driver, raw campaign data and an
independent verifier are available at
\url{https://github.com/kamb-code/sha256-r19-preimage}.
Any of the seven preimages (Table~\ref{tab:preimages}) can be checked with
\begin{center}
\small\texttt{python3 code/verify\_r19.py {-}{-}rounds 19 {-}{-}hash <hex> {-}{-}words "<W0 \ldots W15>"}
\end{center}

%-----------------------------------------------------------------------
\appendix
\section{An Exact Cancellation in the Message Schedule}
\label{app:ninestep}
%-----------------------------------------------------------------------

The $\sigma_0$ analysis above also exposes a two-word message difference
that the schedule cancels exactly for eight consecutive positions.

\begin{lemma}[Nine-step cancellation]
\label{lem:nine-step}
Let $W_0,\ldots,W_{15}$ be any message block and let $\delta\ne0$.  Define
$W'_i = W_i$ for $i\notin\{0,9\}$, $W'_0 = W_0+\delta$ and
$W'_9 = W_9-\delta$ (all arithmetic mod $2^{32}$).  Then the extended
schedules satisfy $W'_r - W_r = 0$ for $r = 16,\ldots,23$, and $W'_{25} - W_{25} = -\delta$, for every block and every $\delta$.  At $r=24$ the term $\sig{0}(W_9)$ breaks the cancellation and $W'_{24}-W_{24} = \sig{0}(W_9-\delta)-\sig{0}(W_9)$ is nonzero in general.
\end{lemma}

\begin{proof}
The schedule recurrence is
$W_r = \sigma_1(W_{r-2}) + W_{r-7} + \sigma_0(W_{r-15}) + W_{r-16}$.
For $r\in[16,23]$ the term $\sigma_0(W_{r-15})$ indexes $r-15\in[1,8]$,
where the difference is zero; the term $\sigma_1(W_{r-2})$ indexes
$r-2\in[14,21]$, where the difference is zero up to $r=23$ by induction;
the term $W_{r-7}$ indexes $r-7\in[9,16]$ and contributes $-\delta$ only
at $r=16$; and the term $W_{r-16}$ indexes $r-16\in[0,7]$ and contributes
$+\delta$ only at $r=16$.  At $r=16$ the two contributions cancel, and for
$r\in[17,23]$ all four terms contribute zero.  At $r=25$ the only nonzero
term is $W_{r-16}=W_9$, giving $-\delta$.
\end{proof}

An exhaustive check over all 240 ordered pairs of positions in $[0,15]$
shows that $(0,9)$ and $(9,0)$ are the only pairs for which $\Delta W_r$
vanishes identically for all $r\in[16,23]$ (script
\texttt{alt\_differential.py}).  The reason is that $\sigma_0(W_k)$
first enters the schedule at $W_{k+15}$, so a perturbed position
$k\in[1,8]$ seeds a nonzero difference inside the window, and cancellation
of the two linear terms at $r=16$ then requires exactly the pair
$\{0,9\}$.

\paragraph{Why no search result is claimed.}
GPU sweeps built on this difference located message pairs with one equal
output word of the full 64-round compression function, and multi-block
coordinate descent reached output differences of Hamming weight $72/256$.
We record these only to state why they are not results.  For a fixed input
difference, the expected minimum Hamming weight of a random 256-bit output difference over $2^{32}$ independent trials is $79$ (median $78$), and over the $2^{37}$--$2^{39}$ evaluations used by the multi-block searches it is $75$--$74$; a minimum of $72$ or below occurs by chance with probability $0.11$--$0.38$ over that many independent trials, and a directed descent does at least as well.  The
observed minima ($75$--$95$ over single $2^{32}$ sweeps, $72$ after
multi-block descent) are therefore at the level expected of an unbiased
difference, and one-word equalities occurred at the uniform rate of $2^{-32}$ per output word per trial.  The code and data remain in the repository for
completeness.

%-----------------------------------------------------------------------
\paragraph{Note added, September 2026.}  After this paper was written, a chosen-context variant of the attack~\cite{Basra2026Submask} imposed three conditions on the context words ($a_4=a_5$, $e_8=e_9=\mathtt{0xFFFFFFFF}$), which removes the fixed-point iteration and collapses the consistency check to a bitwise condition.  It reduces the cost of a 19-round preimage from about $2^{38.3}$ to about $2^{15.3}$ swept $a_0$ and reaches 20 rounds, where the fourth constraint remains an exact $2^{-32}$ filter, giving about $2^{45.4}$ swept $a_0$ per preimage (about 15 hours on one A100); two computed 20-round preimages are reported there, one of the all-ones digest.  That paper supersedes the cost figures here, and the screening result of \S\ref{sec:screening} is of historical interest only for the attack it was measured on.  The figures in this paper are those of the random-context attack and remain correct for it.

\begin{thebibliography}{10}

\bibitem{Basra2026Submask}
K.~S.~Basra.
\newblock Context shaping for reduced-round {SHA-256} compression-function preimages: nineteen rounds in milliseconds, and computed preimages at twenty rounds.
\newblock Manuscript, 2026.  Source, code and witnesses at \url{https://github.com/kamb-code/sha256-r19-preimage}.

\bibitem{Basra2026R20}
K.~S.~Basra.
\newblock The 20-round barrier for table-based preimage attacks on
  {SHA-256}.
\newblock Companion note, 2026.  Source and code in the same repository.

\bibitem{FIPS180}
National Institute of Standards and Technology.
\newblock \emph{{FIPS PUB 180-4}: Secure Hash Standard ({SHS})}, 2015.

\bibitem{IsobeShibutani2009}
T.~Isobe and K.~Shibutani.
\newblock Preimage attacks on reduced {Tiger} and {SHA-2}.
\newblock In \emph{Fast Software Encryption (FSE 2009)}, LNCS 5665,
  pp.~139--155. Springer, 2009.

\bibitem{AokiGuo2009}
K.~Aoki, J.~Guo, K.~Matusiewicz, Y.~Sasaki, and L.~Wang.
\newblock Preimages for step-reduced {SHA-2}.
\newblock In \emph{Advances in Cryptology -- ASIACRYPT 2009}, LNCS 5912,
  pp.~578--597. Springer, 2009.

\bibitem{GuoLiLiuWangZhang2026}
J.~Guo, H.~Li, M.~Liu, S.~Wang, and T.~Zhang.
\newblock Dual-syncopation meet-in-the-middle attacks: new results on
  {SHA-2} and {MD5}.
\newblock In \emph{Advances in Cryptology -- EUROCRYPT 2026, Part VI},
  LNCS 16546, pp.~121--151. Springer, 2026.
  DOI 10.1007/978-3-032-25333-0\_5.  Cryptology ePrint Archive, Report
  2026/353.

\bibitem{KRS2012}
D.~Khovratovich, C.~Rechberger, and A.~Savelieva.
\newblock Bicliques for preimages: attacks on {Skein-512} and the {SHA-2}
  family.
\newblock In \emph{Fast Software Encryption (FSE 2012)}, LNCS 7549,
  pp.~244--263. Springer, 2012.

\bibitem{SanadhyaSarkar2008}
S.~K.~Sanadhya and P.~Sarkar.
\newblock New collision attacks against up to 24-step {SHA-2}.
\newblock In \emph{Progress in Cryptology -- INDOCRYPT 2008}, LNCS 5365,
  pp.~91--103. Springer, 2008.

\bibitem{MendelNadSchlaffer2013}
F.~Mendel, T.~Nad, and M.~Schl\"affer.
\newblock Improving local collisions: new attacks on reduced {SHA-256}.
\newblock In \emph{Advances in Cryptology -- EUROCRYPT 2013}, LNCS 7881,
  pp.~262--278. Springer, 2013.

\bibitem{LiLiuWang2024}
Y.~Li, F.~Liu, and G.~Wang.
\newblock New records in collision attacks on {SHA-2}.
\newblock In \emph{Advances in Cryptology -- EUROCRYPT 2024, Part I}, LNCS 14651,
  pp.~158--186. Springer, 2024.  DOI 10.1007/978-3-031-58716-0\_6.

\bibitem{LiLiuWangDongSun2024}
Y.~Li, F.~Liu, G.~Wang, X.~Dong, and S.~Sun.
\newblock The first practical collision for 31-step {SHA-256}.
\newblock In \emph{Advances in Cryptology -- ASIACRYPT 2024, Part VII}, LNCS 15490,
  pp.~237--266. Springer, 2024.  DOI 10.1007/978-981-96-0941-3\_8.

\bibitem{MironovZhang2006}
I.~Mironov and L.~Zhang.
\newblock Applications of {SAT} solvers to cryptanalysis of hash
  functions.
\newblock In \emph{Theory and Applications of Satisfiability Testing (SAT
  2006)}, LNCS 4121, pp.~102--115. Springer, 2006.

\bibitem{Davydov2025}
V.~V.~Davydov, M.~D.~Pikhtovnikov, A.~P.~Kiryanova, and O.~S.~Zaikin.
\newblock Analysis of the cryptographic strength of the {SHA-256} hash
  function using the {SAT} approach.
\newblock \emph{Scientific and Technical Journal of Information
  Technologies, Mechanics and Optics}, 25(3):428--437, 2025.
  DOI 10.17586/2226-1494-2025-25-3-428-437.

\bibitem{Zaikin2026}
O.~Zaikin.
\newblock Preimage attacks on round-reduced {MD5}, {SHA-1}, and {SHA-256}
  using parameterized {SAT} solver.
\newblock \emph{Constraints}, 31:5, 2026 (received 28 March 2025, published online 23 January 2026).  DOI 10.1007/s10601-025-09383-0.

\end{thebibliography}
\end{document}

